\uuid{HRqF}
\exo7id{3647}
\titre{exo7 3647}
\auteur{quercia}
\organisation{exo7}
\datecreate{2010-03-10}
\isIndication{false}
\isCorrection{false}
\chapitre{Endomorphisme particulier}
\sousChapitre{Autre}
\cpge{mpsi-mp2i.applications-lineaires.e, pcsi.applications-lineaires.f, mpsi-mp2i.dimension-finie.a, pcsi.dimension-finie.a, ptsi.dimension-finie.a, mpsi-mp2i.applications-lineaires.a, pcsi.applications-lineaires.a, ptsi.applications-lineaires.a, bcpst2.applications-lineaires-matrices.b}
\module{Algèbre}
\niveau{L2}
\difficulte{}

\contenu{
\question{
Soit $E$ un $ K$-ev. On suppose qu'il existe $p$ formes linéaires
$f_1,\dots,f_p$ telles que :
$$\forall\ \vec x \in E,\ (f_1(\vec x) = \dots = f_p(\vec x) = 0)
   \Rightarrow  (\vec x = \vec 0).$$
Montrer que $E$ est de dimension finie inférieure ou égale à $p$.
}
}
